\chapter{Acquired Skills During Industrial Training}
\label{chapter-4}

\section{Introduction}

\par Before this internship, My academic knowledge on Javascript and self taught React and NextJS knowledge was basic. I could build functional websites with that but WeCycle changed that. It teaches me the difference between functional and production ready code. It also help me to escape tutorial hell and start utilizing documentations. I learned how to make any UI designs pixel perfect using code. Go beyond my work scope to learn new things and implement backend logic also. WeCycle gave me that opportunity and give a headstart on my career.

\section{Learning New Technologies and Applying Them to Solve
Engineering Problems}

\subsection{React.js and Next.js \cite{reactlearn2024, nextjsdocs2024}}
\begin{itemize}
    \item \textbf{Proficiency Level:} From "Tutorial Hell" to Building Production Features.
    \item \textbf{The Real Skills:}
    \begin{itemize}
        \item \textbf{Component based Architecture:} Dummy project did not care about maintainability, so we can write the whole code in one file. But for production ready code, we need to think long term maintenance. Where modular code and abstraction hides complexity and make it easier to maintain.
        \item \textbf{State Management:} Most of my University projects, I used a lot of state which was not a problem back then. But when comes to large codebases every rerender matters. There code be 100 lines of state which is scary to manage. Where I learn that we can use React context api and hooks to abstract the complexity.
        \item \textbf{Next.js:} Understanding when to use client side rendering and when to use server side rendering was eye opening. I was using client side rendering in every places but When we need SEO and performance over functionality server side rendering is way to go.
    \end{itemize}
\end{itemize}

\subsection{TypeScript}
\begin{itemize}
    \item \textbf{Proficiency Level:} Intermediate.
    \item \textbf{The Real Skills:}
    \begin{itemize}
        \item \textbf{Type Safety:} In my personal projects, I always use keyword "any" for most cases. But I can do it on real world applications. I had to learn the proper way to use typescripts type safety.
        \item \textbf{Generics:} In large codebase, reusability is the key. Generic type is supercharge to make the code reusable.
    \end{itemize}
\end{itemize}

\subsection{Tailwind CSS, and Responsive Design \cite{tailwindcss2024}}
\begin{itemize}
    \item \textbf{Proficiency Level:} Advanced.
    \item \textbf{The Real Skills:}
    \begin{itemize}
        \item \textbf{Pixel Perfect Design:} I have to implement all the design made by our designer in code. These need to be perfect based on style, aligned and functionality.
        \item \textbf{Tailwind CSS:} I primarily used tailwind CSS for styling. Now, I can implement designs which confidence and it takes less time to implement a design.
    \end{itemize}
\end{itemize}

\subsection{Backend Development}
\begin{itemize}
    \item \textbf{Proficiency Level:} Intermediate
    \item \textbf{The Real Skills:}
    \begin{itemize}
        \item \textbf{Node.js \& Express:} Built and maintained RESTful APIs for the dashboard and mobile app.
        \item \textbf{Database Management:} Worked with MongoDB (NoSQL) and PostgreSQL to manage user data, bookings, and inventory using Prisma/Mongoose.
        \item \textbf{Payment Integration:} Implemented Stripe Connect for splitting payments between platform and haulers.
    \end{itemize}
\end{itemize}

\subsection{DevOps \& Infrastructure}
\begin{itemize}
    \item \textbf{Proficiency Level:} Intermediate
    \item \textbf{The Real Skills:}
    \begin{itemize}
        \item \textbf{Docker:} Containerized the Chatbot service for consistent deployment across environments.
        \item \textbf{Cloud Services:} Managed DigitalOcean droplets, configured Nginx reverse proxies, and set up SSL certificates using Certbot.
        \item \textbf{CI/CD:} Optimized GitHub Actions pipelines to reduce build times.
    \end{itemize}
\end{itemize}

\subsection{Git and Version Control}
\begin{itemize}
    \item \textbf{Proficiency Level:}  Expert
    \item \textbf{The Real Skills:}
    \begin{itemize}
        \item \textbf{Pull Requests:} In my project, I would directly push my code on the main branch. But in real project it is not feasible as the code can break the production. So I learned to use Pull Requests and code review before merging into main.
        \item \textbf{Merge Conflicts:} As different developers are working on different features. Most of the time there was merge conflicts. So, I need to understand the process and resolve the merge conflicts get both features working at the same time.
\end{itemize}


\section{Theoretical Knowledge Gained}

\subsection{Software Development Concepts}
\begin{itemize}
    \item \textbf{Components:} splitting files and modularity is important for large codebases. If I change the button style in one place, it should update everywhere.
    \item \textbf{APIs:} Frontend is useless without data. I learned how to fetch, and handle errors when the backend is not responding as expected.
    \item \textbf{Caching:} Fetching data every time is a bad idea for most of the time. React Query helps browser level caching for that.
\end{itemize}

\subsection{Development Methodologies}
\begin{itemize}
    \item \textbf{Agile:} Daily standups on my work updates and bi-weekly sprint planning forces me to align with business goals.
    \item \textbf{Code Reviews:} The most humbling part. I get to know that I do not know everything. There is a lot to learn. Code review is not insulting, it is learning opportunity.
\end{itemize}

\section{Practical Knowledge Acquired}

\par The following things I have learnt from WeCycle.

\begin{itemize}
    \item \textbf{The Worksflow:}
    \begin{itemize}
        \item Don't reinvent the wheel. Use open source tools and library so speed up the process.
        \item Use latest but reliable tools and update tools regularly which make the developer experience way better.
    \end{itemize}
    \item \textbf{Debugging:}
    \begin{itemize}
        \item Learn to use React DevTools, it is much better than traditional console logging every variable.
        \item Use Postman and other tools which are helpful for api related bugs.
    \end{itemize}
    \item \textbf{Startup Logistics:}
    \begin{itemize}
        \item Build a simple team which designer, frontend and backend developer and create MVP as fast as possible. Then as the product is start generating revenue, the team can be expanded.
    \end{itemize}
\end{itemize}

\section{Communication Skills Developed}

\begin{itemize}
    \item \textbf{Communication:}
    \begin{itemize}
        \item \textbf{Writing Tickets:} I need to write my own Jira tickets for my own task.
        \item \textbf{Asking Questions:} I stop assuming and start asking questions which helped me to gain additional information and reduce my after work refinement.
    \end{itemize}
    \item \textbf{Teamwork:}
    \begin{itemize}
        \item \textbf{Standups:} I learn to explain my whole day of work in short 5 minutes meetings.
        \item \textbf{Remote Etiquette:} I try to keep my camera open for most of the time which help me connect better.
    \end{itemize}
\end{itemize}

\section{Collaborating Without Killing Each Other}

\begin{itemize}
    \item \textbf{Designers:} They want pixel perfection. So I need to make sure everything is according to the design. Sometime, I don't find something which should be on the design instead of blaming I simplly ask for changes.
    \item \textbf{Backend:} They need the proper utilization of the server. So, I need to fetch data when needed and utilize caching and other technics to minimize the load from the server.
\end{itemize}

\section{Ethics \& Professionalism}

\par In professional environment trust is everything. As we are working remotely there is no way to monitor every single steps we make during our working hours. We do not lie about our work with others. We try to write clean code. Writing bad code is a liability for the organization.

\begin{table}[H]
\centering
\caption{How I Kept It Professional}
\begin{adjustbox}{width=0.9\textwidth}
\begin{tabular}{|l|p{10cm}|}
\hline
\textbf{Principle} & \textbf{What I Actually Did} \\ \hline
\textbf{User Safety} & I made sure buttons were big enough to tap and text was readable. Bad UI frustrates users. \\ \hline
\textbf{Honesty} & If I was going to miss a deadline, I said so. Hiding it makes it worse. \\ \hline
\textbf{Inclusion} & I tested on cheap Android phones, not just my expensive MacBook. Everyone deserves a working app. \\ \hline
\textbf{Competence} & I admitted when I didn't know something and then went and learned it. \\ \hline
\end{tabular}
\end{adjustbox}
\label{tab:ieee-ethics-mapping}
\end{table}

\section{Conclusion}

\par I have successfully completed my training and learnt a lot of things during the process. I learnt how to make the scalable and production ready applications. My proficiency on the tools and technologies improved and get to use unfamiliar tools which are not familiar to me.